% Technical supplement reconstructed from the authors' SIGMOD appendix.
\documentclass[conference]{IEEEtran}
\usepackage{amsmath,amssymb,amsfonts}
\usepackage{array,booktabs,tabularx}
\usepackage{graphicx}
\usepackage{capt-of}
\usepackage{xr-hyper}
\usepackage[hidelinks]{hyperref}
\externaldocument[main-]{main}[main.pdf]
\renewcommand{\thetable}{S\arabic{table}}
\newcommand{\col}[1]{\texttt{#1}}
\begin{document}
\title{TAILOR: Technical Supplement}
% \author{\IEEEauthorblockN{Anonymous Author(s)}}
\author{
\IEEEauthorblockN{1\textsuperscript{st} Pengjie Cai}
\IEEEauthorblockA{\textit{Tianjin University}, 
Tianjin, China \\
pengjiecai\_2025@tju.edu.cn}
\and
\IEEEauthorblockN{2\textsuperscript{nd} Haitao Yuan}
\IEEEauthorblockA{\textit{Nanyang Technological University},
Singapore \\
haitao.yuan@ntu.edu.sg}
}
\maketitle
\appendices

\section{Full Ablation Results}
\label{supp:ablation-section}
Table~\ref{supp:ablation} extends main Table~\ref{main-tab:ablation} to Adult, Default, and Magic. Variant definitions and statistical aggregation follow main Sections~\ref{main-sec:ablation} and~\ref{main-sec:setup}, respectively.

\par\medskip\noindent\begin{minipage}{\linewidth}
\centering
\captionof{table}{Full ablations on the three additional datasets.}
\label{supp:ablation}
\includegraphics[width=\linewidth]{figures/TABLE-3-REST.pdf}
\end{minipage}\par\medskip

On these three datasets, adding the floor raises MOS over Select, and repair raises it further. Full TAILOR has the highest MOS, while the intermediate search variants are not consistently ordered. As in the main ablation, the +MCTS/+LLM comparison includes changes to both initialization and subsequent proposals.

\section{Proofs}
\label{supp:proofs}
We prove main Propositions~\ref{main-prop:operator} and~\ref{main-prop:scalar}, Lemma~\ref{main-lem:metric}, and the search bounds. All algorithm references refer to the main manuscript. The operator proofs assume fixed processing references and row indicators.

\subsection{Dominance Consistency}
\begin{IEEEproof}[Proof of Proposition~\ref{main-prop:operator}(i)]
The relation $\succ$ is irreflexive: $a\succ a$ would require one component to exceed itself. It is transitive because componentwise inequalities compose, and a strict gain on at least one component survives. Consequently it has no cycle.

For a row $x$, let $\ell(x)$ be the largest number of nodes in a domination chain ending at $x$, with $\ell(x)=1$ if $x$ has no dominator. Non-dominated sorting assigns $\operatorname{front}(x)=\ell(x)$. To see this, the first front consists exactly of nodes without a dominator. Inductively, after removing fronts $1,\ldots,f-1$, a row enters front $f$ exactly when all its dominators are in earlier fronts and some dominator is in front $f-1$. Its longest predecessor chain therefore has $f-1$ nodes, and appending the row gives a chain of length $f$.

If $a\succ b$, a longest chain ending at $a$ can be extended by $b$. Acyclicity ensures that no node repeats. Hence $\ell(b)\ge\ell(a)+1$, which yields $\operatorname{front}(b)>\operatorname{front}(a)$.
\end{IEEEproof}

\subsection{Reference-Based Fidelity Floors}
\begin{IEEEproof}[Proof of Proposition~\ref{main-prop:operator}(ii)]
The thresholds in Algorithm~\ref{main-alg:postprocessing}, line~\ref{main-alg:line:tau}, are fixed as $\tau_{\mathrm{fid}}=\beta G_{\mathrm{fid}}(A_{\mathrm{ref}})$ and $\tau_{\mathrm{trend}}=\beta G_{\mathrm{trend}}(A_{\mathrm{ref}})$. The output gate has two branches. If the candidate passes, then $A^\star=A$ and the tests give exactly the required inequalities, together with $|A|=k$. Otherwise $A^\star=A_{\mathrm{ref}}$. Since $0<\beta\le1$ and fidelity scores are nonnegative,
\[
G_q(A_{\mathrm{ref}})\ge\beta G_q(A_{\mathrm{ref}})=\tau_q,
\qquad q\in\{\mathrm{fid},\mathrm{trend}\}.
\]
Both branches therefore satisfy the floors. The reference also provides a size-$k$ fallback, assuming that the feasible preselected pool contains at least $k$ rows.

The thresholds are fractions of fidelity scores achieved by an actual candidate subset. Arbitrary external thresholds could exceed the fidelity attainable within the pool, invalidating the fallback argument. The guarantee compares the release with its constructed reference and does not imply improvement over Raw in every metric.
\end{IEEEproof}

\subsection{Repair Descent and Residual Deviation}
\begin{IEEEproof}[Proof of Proposition~\ref{main-prop:operator}(iii)]
The pool is finite and the subset size is fixed, so only finitely many subsets and deviation values $\phi(A)=|G_{\mathrm{dcr}}(A)-0.5|$ are available. Accepted swaps strictly decrease $\phi$; a sequence of such swaps cannot revisit a subset and has finite length. Fidelity and utility checks only restrict the admissible neighborhood.

Let $h=m/k$. An accepted batch exchanges $m$ rows with opposite balance indicators, so processing DCR changes by exactly $h$ toward the target. A non-crossing move reduces $\phi$ by $h$. A crossing move is improving only if $\phi_{\mathrm{old}}>h/2$ and leaves
\[
\phi_{\mathrm{new}}=h-\phi_{\mathrm{old}}<h/2.
\]
From this residual, another move of size $h$ in either direction cannot improve the absolute deviation. There are at most $\lfloor\phi(A_0)/h\rfloor$ non-crossing accepted swaps and at most one crossing swap. If the initial ratio is an integer, the non-crossing sequence reaches zero and no crossing swap is accepted. Hence the bound is
\[
\#\text{accepted swaps}\le\left\lceil\frac{k\phi(A_0)}m\right\rceil.
\]
A finite repair budget may stop the procedure earlier and cannot increase this count.

Suppose now that the budget does not bind and an admissible swap preserving fidelity and utility exists whenever $\phi>h/2$. A move by $h$ toward $0.5$ would then strictly reduce $\phi$ and be accepted. Consequently termination can occur only at $\phi(A_{\mathrm{rep}})\le h/2=m/(2k)$. This deviation bound transfers to the final output when the gate releases the repaired subset; the fallback reference is not covered by it.

The pre-repair subset satisfies the floors, and every accepted swap preserves them. Induction on $\bar U(A_{t+1})\ge\bar U(A_t)$ also shows that mean surrogate utility never drops below its pre-repair value. This bound applies to the row score; test performance of a retrained predictor can still change.
\end{IEEEproof}

\subsection{Scalarization Properties}
\begin{IEEEproof}[Proof of Proposition~\ref{main-prop:scalar}]
Write $\kappa=\rho/(1+\rho)$, $\tilde z_j=(z_j+\rho)/(1+\rho)$, and $\Gamma=\prod_j\tilde z_j^{w_j}$. Because $z_j\in[0,1]$, each $\tilde z_j\in[\kappa,1]$. Positive weights summing to one imply $\kappa\le\Gamma\le1$. The reward can be written as
\[
J=\frac{\Gamma-\kappa}{1-\kappa}=(1+\rho)\Gamma-\rho\in[0,1].
\]

\textit{(i) Arithmetic-mean bound.} Weighted AM--GM gives $\Gamma\le\sum_jw_j\tilde z_j$, with equality iff all $\tilde z_j$ are equal. Applying the increasing affine transformation above and using $\sum_jw_j=1$ gives
\[
J\le(1+\rho)\sum_jw_j\tilde z_j-\rho=\sum_jw_jz_j.
\]
The transformation from $z_j$ to $\tilde z_j$ is injective, so equality holds iff all original components are equal.

\textit{(ii) Strict monotonicity.} For each $j$,
\[
\frac{\partial\Gamma}{\partial z_j}
=\frac{w_j}{1+\rho}\tilde z_j^{w_j-1}
\prod_{q\ne j}\tilde z_q^{w_q}>0.
\]
All factors are positive, including at boundary values because $\kappa>0$. Since $J$ is an increasing affine function of $\Gamma$, it is strictly increasing in each component throughout the unit cube.

\textit{(iii) Single-failure ceiling.} Fix $z_j=0$, hence $\tilde z_j=\kappa$. Strict monotonicity makes the maximum occur at $z_q=1$ for every $q\ne j$. Therefore
\[
\max_{z_j=0}J=\frac{\kappa^{w_j}-\kappa}{1-\kappa}<1.
\]
With $\rho=0.05$ and $w_j=1/4$, the ceiling is approximately $0.44$, even when every other component equals one.

\textit{(iv) Pareto consistency.} Moving from a dominated vector $z(A')$ to its dominating vector $z(A)$ increases at least one coordinate and decreases none. Part (ii) gives $J(A)>J(A')$.

\textit{(v) Optimality within the reachable family.} If a reachable subset $A(\Theta')\in\mathcal R$ Pareto-dominated an exact maximizer $A(\Theta^\star)$, part (iv) would imply a strictly larger reward, contradicting maximality. The family $\mathcal R$ is limited by the operator, policy validity, and any horizon restriction. This proof does not establish that finite MCTS evaluates the maximizer.
\end{IEEEproof}

\subsection{Metric-Dominant Exploitation}
\begin{IEEEproof}[Proof of Lemma~\ref{main-lem:metric}]
Both confidence values lie in $[0,1]$, so their contribution to the base score differs by at most $0.05$. Thus
\[
\begin{split}
B(v_1)-B(v_2)
&\ge R^\star(v_1)-R^\star(v_2)-0.05\\
&>0.
\end{split}
\]
The conclusion concerns $B$. The exploration term can still change the UCT ranking, and smaller reward gaps need not preserve the base-score ordering.
\end{IEEEproof}

\subsection{Search Value Invariant and Bounds}
\label{supp:search-bounds}
\begin{IEEEproof}[Proof of the anytime bound]
After root evaluation, the root's stored reward equals its measured value. Assume each stored $R^\star(v)$ equals the best measured reward in the subtree rooted at $v$. A new rollout adds a measured reward $R_i$ to one node's subtree and those of its ancestors. Algorithm~\ref{main-alg:mcts}, line~\ref{main-alg2:line:backup}, updates their stored values to the maximum of the previous value and $R_i$. All other subtrees remain unchanged, preserving the invariant after each evaluation.

The evaluated policy set grows and never loses a member. Its maximum is nondecreasing and is at least the initialized root reward. The returned policy attains that maximum, proving the main bound. All candidate releases execute Algorithm~\ref{main-alg:postprocessing}, so they inherit its fidelity gate. These statements apply to the same pool, references, and reward used in search.
\end{IEEEproof}

\begin{IEEEproof}[Proof of the conditional convergence bound]
Assume the specified improving sequence has length $L$ and, conditional on the full history, each round advances it with probability at least $p$. Let $E_{\mathrm{adv}}$ count advances among $E$ rounds, stopping the count once the target is reached. Before success, the conditional lower bound permits a coupling with independent Bernoulli trials of success probability $p$. The chance of fewer than $L$ advances is therefore bounded by the lower tail of $\operatorname{Bin}(E,p)$.

For $E\ge\max\{2L/p,8\ln(1/\delta)/p\}$, we have $L\le pE/2$. A multiplicative Chernoff bound gives
\[
\begin{split}
\Pr[E_{\mathrm{adv}}<L]
&\le\Pr[\operatorname{Bin}(E,p)<pE/2]\\
&\le\exp(-pE/8)\le\delta.
\end{split}
\]
Once the optimum is evaluated, its reward is retained by the invariant above. The bound requires the stated sequence-advancement assumption, which is not established by UCT alone.
\end{IEEEproof}

\onecolumn
\section{Compressed Prompt Specification}
\label{supp:prompts}
\subsection{Compression Motivation}
The LLM receives structured summaries of dataset semantics, source diagnostics, and rollout feedback, excluding full candidate tables and complete experiment traces. It proposes column policies and, in the auxiliary Mixed setting, source composition. The deterministic operator selects rows, and measured quality supplies the reward.

\subsection{Shared Context Blocks}
Table~\ref{supp:prompt-blocks} lists the context fields used to assemble prompts, combining fixed dataset/source summaries with diagnostics updated during search. It specifies the interface contents; the entries are not verbatim prompts.

\par\medskip\noindent\begin{minipage}{\textwidth}
\centering
\captionof{table}{Compressed context blocks used by TAILOR.}
\label{supp:prompt-blocks}
\small
\begin{tabularx}{\textwidth}{@{}p{0.16\textwidth}X p{0.26\textwidth}@{}}
\toprule
Block & Content & Purpose\\
\midrule
Dataset brief & Dataset name, target, column types, priority columns, target summary & Column semantics and task background\\
Theta size bounds & Minimum and maximum sizes of $C_{1D},C_{2D},C_P$, together with target rules & Admissible policy constraints\\
DCR semantics & Balance target near $0.5$ and reward $1-|\mathrm{DCR}-0.5|$ & Avoid treating DCR as a monotonic objective\\
Real utility profile & Split sizes, evaluator, metric, leading utility features, semantic summary & Guide $c_U$ and task-sensitive edits\\
Source profile & Mean metrics, weak marginals and pairs, DCR imbalance, utility features & Source diagnosis; composition edits in Mixed\\
Rollout feedback & Evaluated policies, rewards, weak columns/pairs, DCR and source contribution & Node diagnosis and refinement\\
Reference nodes & Best measured-reward node, best LLM-score node, siblings, global best & Context and duplicate avoidance\\
\bottomrule
\end{tabularx}
\end{minipage}\par\medskip

\subsection{Prompt Interfaces and Validation}
Single-source search uses three interfaces: initialization constructs a valid policy, diagnosis summarizes limiting objectives, and refinement proposes targeted edits. Mixed-source search also summarizes real and synthetic sources and supports hybrid-pool initialization and refinement. Validation checks schema membership, component-size bounds, target rules, duplicate policies, and source multipliers where applicable. Invalid proposals are repaired or replaced by a deterministic fallback. Each admitted policy then executes the operator and receives a measured reward.

\section{Adult Feedback-to-Policy Interpretability Case}
\label{supp:adult}
Tables~\ref{supp:adult-feedback} and~\ref{supp:adult-policy} show the recorded Adult diagnostics and policy components. The trace links measured weaknesses to proposed changes in the columns used by each objective.

\par\medskip\noindent\begin{minipage}{\textwidth}
\centering
\captionof{table}{Adult feedback-to-$\Theta$ interpretability case.}
\label{supp:adult-feedback}
\small
\begin{tabularx}{\textwidth}{@{}p{0.13\textwidth}X X X@{}}
\toprule
Signal & Recorded evidence & Policy response & Interpretation\\
\midrule
Worst shape & \col{education.num} appears among weak marginal columns & Add \col{education.num} to $C_{1D}$ & Focus marginal scoring on a diagnosed weakness\\
Worst trend & The \col{marital.status}--\col{sex} pair appears among weak dependencies & Include the diagnosed pair in $C_{2D}$ & Convert pairwise feedback into a trend focus\\
Utility features & \col{marital.status} and \col{occupation} rank highly & Replace or adjust $c_U$ & Change the feature defining utility groups\\
DCR quantiles & Balance is near its target but the nearest-neighbor representation requires stability & Expand $C_P$ & Strengthen the feature space used by the distance surrogate\\
\bottomrule
\end{tabularx}
\end{minipage}\par\medskip

\par\medskip\noindent\begin{minipage}{\textwidth}
\centering
\captionof{table}{Recorded Adult policy components. The trend row lists anchor columns appearing in the recorded pairwise configuration.}
\label{supp:adult-policy}
\small
\begin{tabularx}{\textwidth}{@{}p{0.10\textwidth}p{0.25\textwidth}X@{}}
\toprule
Component & Role & Columns\\
\midrule
$C_{1D}$ & Marginal focus & age, capital.loss, education.num, fnlwgt, income, marital.status, native.country, occupation, race, sex, workclass\\
$C_{2D}$ & Pairwise trend anchors & age, education, education.num, fnlwgt, hours.per.week, marital.status, occupation, race, relationship, sex, workclass\\
$C_P$ & DCR-aware features & age, capital.gain, capital.loss, education, education.num, fnlwgt, hours.per.week, marital.status, native.country, occupation, race, relationship, sex, workclass\\
$c_U$ & Utility balance anchor & capital.gain\\
\bottomrule
\end{tabularx}
\end{minipage}\par\medskip

The $C_{2D}$ row lists participating anchor columns; it does not specify the complete pair set or their Cartesian product. Utility-feature diagnostics suggest candidates for $c_U$, and measured rewards determine the retained anchor. The trace records these choices without isolating the performance contribution of any single column.

\par\medskip\noindent\begin{minipage}[t]{0.48\textwidth}
\section{Additional Privacy Results}
\label{supp:privacy-section}
Table~\ref{supp:privacy} extends main Table~\ref{main-tab:audit} using the same attack data roles and aggregation protocol. Floor and Repair are the operator variants defined in the main ablation; TAILOR also includes learned policy search. Their comparison measures the incremental changes after repair and policy learning.

\par\medskip\noindent\begin{minipage}{\linewidth}
\centering
\captionof{table}{Additional membership-inference results.}
\label{supp:privacy}
\includegraphics[width=\linewidth]{figures/TABLE-4-REST.pdf}
\end{minipage}\par\medskip

\end{minipage}\hfill\begin{minipage}[t]{0.48\textwidth}
\section{Million-Scale Online Execution Profile}
\label{supp:profile-section}
Table~\ref{supp:profile} gives the stage breakdown for one online application of the learned policy in the main million-scale experiment. The total excludes final metric evaluation, as specified in the main timing protocol.

\par\medskip\noindent\begin{minipage}{\linewidth}
\centering
\captionof{table}{Online execution on US Census 1990, in seconds.}
\label{supp:profile}
\begin{tabular}{@{}lr@{}}
\toprule
Stage & Time (s)\\
\midrule
Preselection & $770.49\pm6.32$\\
Scoring & $96.26\pm2.10$\\
Core selection & $28.60\pm0.42$\\
Repair & $215.32\pm1.45$\\
\midrule
Online total & $1110.67\pm8.71$\\
\bottomrule
\end{tabular}
\end{minipage}\par\medskip

Preselection and repair account for approximately $69.4\%$ and $19.4\%$ of online execution time, respectively, making preselection the largest cost in this case.
\end{minipage}\par
\end{document}